\subsection{Raw, Uncompressed Storage Favours Native UTM}
Both representations are fixed in size, so the ratio in
Table~\ref{tab:storage_accounting} is a geometric constant rather than a
measurement that varies: over the 40 regions its range at
$\gamma_{\mathrm{tile}}=1$ is $2.584$ to $2.584$.

\begin{table}[t]
\centering
\caption{Raw storage for the same nominal ground area}
\label{tab:storage_accounting}
\begin{tabular}{lrr}
\toprule
Representation & Raw size (MB) & HEALPix / this \\
\midrule
HEALPix tile (level 20, $1024^2$ cells) & 4.194 & 1.000 \\
Equal-area UTM crop ($637^2$ px) & 1.623 & 2.584 \\
UTM crop $\times\gamma_{\mathrm{tile}}=1.3$ & 2.110 & 1.988 \\
\bottomrule
\end{tabular}
\end{table}


Table~\ref{tab:storage_accounting} reports the raw storage accounting for the two grids. Because both the HEALPix level-20 tile and the equal-area UTM crop of Section~\ref{sec:utm_crop} are fixed-size representations of the same nominal ground area, their raw byte counts are geometric constants, identical for every scene. At $\gamma_{\mathrm{tile}}=1$ (no assumed overhead) the HEALPix tile is $2.584\times$ larger than the UTM crop, which is exactly $1024^2/637^2$ and therefore identical for every region by construction; it matches the analytic oversampling factor $\alpha_{\mathrm{HPX/native}}\approx2.58$ of \eqref{eq:alpha} because the crop side was derived from it. This is the geometry-only result, and the only one in this paper that is exact. Under the illustrative $\gamma_{\mathrm{tile}}=1.30$ Sentinel-2/MGRS granule-overhead scenario the gap narrows but does not close: the uncompressed HEALPix tile remains $1.99\times$ the size of the UTM crop scaled by $\gamma_{\mathrm{tile}}$. Raw storage alone therefore does not make HEALPix competitive under either scenario, since the $2.58\times$ geometric oversampling of level 20 is too large for a $30\%$ overhead to offset. Whether HEALPix can nonetheless become competitive is a question that only compression, by exploiting the redundancy this oversampling represents, can answer.
